% Part: first-order-logic
% Chapter: axiomatic-deduction
% Section: proving-things

\documentclass[../../../include/open-logic-section]{subfiles}

\begin{document}

\iftag{FOL}
      {\olfileid[id]{fol}{axd}{pro}}
      {\olfileid[id]{pl}{axd}{pro}}

\olsection{Contoh \usetoken{P}{derivation}}


\begin{ex}
Misalkan kita ingin membuktikan $(\lnot !D \lor !E) \lif (!D \lif
!E)$. Jelas bahwa ini bukan instansiasi dari aksioma mana pun, sehingga
kita harus menggunakan aturan \MP{} untuk !!{derive}nya. Aturan yang
kita perlukan di sini adalah~MP, yang, jika diberikan $!A$ dan $!A \lif !B$,
memungkinkan kita membenarkan~$!B$. Salah satu strategi ialah menggunakan
\olref[prp]{ax:lor3} dengan $!A$ sebagai $\lnot !D$, $!B$ sebagai $!E$,
dan $!C$ sebagai $!D \lif !E$, yakni instansiasi
\[
(\lnot !D \lif (!D \lif !E)) \lif ((!E \lif (!D \lif !E)) \lif ((\lnot
!D \lor !E) \lif (!D \lif !E))).
\]
Mengapa? Dua penerapan MP menghasilkan bagian terakhir, dan bagian itulah
yang kita inginkan. Kita juga mudah melihat bahwa $\lnot !D \lif
(!D \lif !E)$ merupakan instansiasi dari \olref[prp]{ax:lnot2}, dan
$!E \lif (!D \lif !E)$ merupakan instansiasi dari
\olref[prp]{ax:lif1}. Jadi, !!{derivation} kita adalah:
\begin{derivation}
  1. &  $\lnot !D \lif (!D \lif !E)$ & \olref[prp]{ax:lnot2} \\
  2. & $(\lnot !D \lif (!D \lif !E)) \lif {}$\\
  &\qquad $((!E \lif (!D \lif !E)) \lif ((\lnot !D \lor !E) \lif (!D \lif !E)))$ & \olref[prp]{ax:lor3}\\
  3. & $(!E \lif (!D \lif !E)) \lif ((\lnot !D \lor !E) \lif (!D \lif !E))$ & 1, 2, \MP\\
  4. & $!E \lif (!D \lif !E)$ & \olref[prp]{ax:lif1}\\
  5. & $(\lnot !D \lor !E) \lif (!D \lif !E)$ & 3, 4, \MP
\end{derivation}
\end{ex}

\begin{ex}\ollabel{ex:identity}
Mari kita coba mencari !!a{derivation} atas $!D \lif !D$. Formula ini
bukan instansiasi suatu aksioma, sehingga kita harus menggunakan \MP{}
untuk !!{derive}nya. \olref[prp]{ax:lif1} adalah aksioma berbentuk
$!A \lif !B$ yang dapat kita kenai~\MP. Agar berguna, tentu saja $!B$
yang dibenarkan sebagai langkah benar oleh \MP{} dalam kasus ini haruslah
$!D \lif !D$, sebab itulah yang ingin kita !!{derive}. Ini berarti
$!A$ juga harus berupa $!D$; dengan kata lain, kita dapat memeriksa
instansiasi \olref[prp]{ax:lif1} berikut:
\[
!D \lif (!D \lif !D)
\]
Untuk menerapkan \MP, kita juga harus membenarkan premis kedua yang
bersesuaian, yakni~$!A$. Namun, dalam kasus kita premis itu adalah~$!D$,
dan kita tidak akan dapat !!{derive}~$!D$ dengan sendirinya. Jadi, kita
memerlukan strategi lain.

Aksioma lain yang hanya melibatkan~$\lif$ adalah
\olref[prp]{ax:lif2}, yakni
\[
(!A \lif (!B \lif !C)) \lif ((!A \lif !B) \lif (!A \lif !C))
\]
Kita dapat mencapai implikasi bersarang terakhir dengan menerapkan
\MP{} dua kali. Sekali lagi, ini berarti kita menginginkan instansiasi
\olref[prp]{ax:lif2} dengan $!A \lif !C$ berupa $!D \lif !D$, yaitu
!!{formula} yang kita tuju. Dengan demikian, $!A$ dan $!C$ tentu
keduanya adalah~$!D$. Bagaimana sebaiknya kita memilih~$!B$ agar baik
$!A \lif (!B \lif !C)$ maupun $!A \lif !B$, yakni dalam kasus kita
$!D \lif (!B \lif !D)$ dan $!D \lif !B$, juga !!{derivable}? Formula
pertama sudah merupakan instansiasi dari \olref[prp]{ax:lif1}, apa pun
pilihan kita untuk~$!B$. Sementara itu, $!D \lif !B$ akan menjadi
instansiasi lain dari \olref[prp]{ax:lif1} jika $!B$ berupa
$(!D \lif !D)$. Jadi, !!{derivation} kita adalah:
\begin{derivation}
  1. & $!D \lif ((!D \lif !D) \lif !D)$ & \olref[prp]{ax:lif1}\\
  2. & $(!D \lif ((!D \lif !D) \lif !D)) \lif {}$\\
  & \qquad $((!D \lif (!D \lif !D)) \lif (!D \lif !D))$ & \olref[prp]{ax:lif2}\\
  3. & $(!D \lif (!D \lif !D)) \lif (!D \lif !D)$ & 1, 2, \MP\\
  4. & $!D \lif (!D \lif !D)$ & \olref[prp]{ax:lif1}\\
  5. & $!D \lif !D$ & 3, 4, \MP
\end{derivation}
\end{ex}

\begin{ex}\ollabel{ex:chain}
Kadang-kadang kita ingin menunjukkan bahwa terdapat !!a{derivation}
atas suatu !!{formula} dari sejumlah !!{formula} lain~$\Gamma$. Sebagai
contoh, mari kita tunjukkan bahwa kita dapat !!{derive} $!A \lif !C$
dari $\Gamma = \{!A \lif !B, !B \lif !C\}$.
\begin{derivation}
  1. & $!A \lif !B$ & \Hyp\\
  2. & $!B \lif !C$ & \Hyp\\
  3. & $(!B \lif !C) \lif (!A \lif (!B \lif !C))$ & \olref[prp]{ax:lif1} \\
  4. & $!A \lif (!B \lif !C)$ & 2, 3, \MP\\
  5. & $(!A \lif (!B \lif !C)) \lif {}$\\
  & \qquad $((!A \lif !B) \lif (!A \lif !C))$ & \olref[prp]{ax:lif2}\\
  6. & $((!A \lif !B) \lif (!A \lif !C))$ & 4, 5, \MP\\
  7. & $!A \lif !C$ & 1, 6, \MP
\end{derivation}
Baris-baris yang berlabel ``\Hyp'' (untuk ``hipotesis'') menunjukkan
bahwa !!{formula} pada baris tersebut merupakan !!a{element}
dari~$\Gamma$.
\end{ex}

\begin{prop}
  \ollabel{prop:chain} Jika $\Gamma \Proves !A \lif !B$ dan $\Gamma
  \Proves !B \lif !C$, maka $\Gamma \Proves !A \lif !C$.
\end{prop}

\begin{proof}
  Misalkan $\Gamma \Proves !A \lif !B$ dan $\Gamma \Proves !B \lif
  !C$. Maka terdapat !!a{derivation} atas $!A \lif !B$ dari~$\Gamma$,
  serta !!a{derivation} atas~$!B \lif !C$ dari~$\Gamma$. Gabungkan
  keduanya menjadi satu !!{derivation} dengan mengonkatenasikannya.
  Sekarang tambahkan baris 3--7 dari !!{derivation} dalam contoh
  sebelumnya. Hasilnya merupakan !!a{derivation} atas $!A \lif !C$---
  yaitu baris terakhir !!{derivation} baru tersebut---dari~$\Gamma$.
  Perhatikan bahwa pembenaran baris 4 dan~7 tetap berlaku jika acuan
  kepada nomor baris~1 diganti dengan acuan kepada baris terakhir
  !!{derivation} atas~$!A \lif !B$, dan acuan kepada nomor baris~2
  diganti dengan acuan kepada baris terakhir !!{derivation}
  atas~$!B \lif !C$.
\end{proof}

\begin{prob}
Tunjukkan bahwa pernyataan-pernyataan berikut berlaku dengan menyajikan
!!{derivation}s dari aksioma-aksioma:
\begin{enumerate}
\item $(!A \land !B) \lif (!B \land !A)$
\item $((!A \land !B) \lif !C) \lif (!A \lif (!B \lif !C))$
\item $\lnot(!A \lor !B) \lif \lnot !A$
\end{enumerate}
\end{prob}



\end{document}
